\documentclass[../../../main.tex]{subfiles}

\begin{document}

\section{Generalizations}

So far, we focused on functions of two variables. But what about
functions of $n$ variables? The whole point of multivariable calculus
is to generalize results from single variable calculus to functions
of two or more variables. So let's do that! Since we generalized
for functions of two variables, it shouldn't be too difficult
generalizing to three or more variables.

\begin{definition}{}{}
Given a function $f$ of $n$ variables, assume that $f$ is defined
on an open set $D$ that contains $(a_1, \ldots, a_n)$ where all first
partials of $f$ exist at $(a_1, \ldots, a_n)$. Moreover,
let $\Delta f = f(x_1, \ldots, x_n) - f(a_1, \ldots, a_n)$ and
$\Delta x_i = x_i - a_i$. The function $f$ is \textbf{differentiable}
at the point $(a_1, \ldots, a_n)$ if there exists some functions
$\epsilon_1, \ldots, \epsilon_n$ such that
\begin{align*}
    \Delta f
    &= f_{x_1}(a_1, \ldots, a_n) \Delta x_1 + \cdots
        + f_{x_n}(a_1, \ldots, a_n) \Delta x_n \\
    &\qquad\qquad\ + \epsilon_1(x_1, \ldots, x_n) \Delta x_1
        + \cdots + \epsilon_n(x_1, \ldots, x_n) \Delta x_n
\end{align*}
holds for all points in the open set $D$ and $\epsilon_i$ approaches
zero as $(\Delta x_1, \ldots, \Delta x_n) \to (0, \ldots, 0)$.
\end{definition}

Notice that this aligns with Definition \ref{def:Differentiability
for Multivariable} for $n = 2$. Similarly, we can generalize
Theorem \ref{theo:Differentiability and Partial Continuous}.

\begin{theorem}{}{}
Let $f$ be a function of $n$ variables whose first-order partials
are defined on an open ball around $(a_1, \ldots, a_n)$. The function
$f$ is differentiable at $(a_1, \ldots, a_n)$ if the partials are
continuous at $(a_1, \ldots, a_n)$.
\end{theorem}

Similarly, differentiability implies continuity. Now we get more
interesting result for generalizing chain rule. Like this is so
important to machine learning and our lives. If I am getting a
tattoo one day, I will be getting one for this chain rule.

\begin{theorem}{}{}
Let $f$ be a function of $n$ variables that are differentiable
functions of $t$. Then, $f(x_1(t), \ldots, x_n(t))$ is also
differentiable and the following chain rule holds.
\[
\frac{d}{dt} f(x_1(t), \ldots, x_n(t))
= \sum_{i=1}^n \frac{\partial f}{\partial x_i} \frac{dx_i}{dt}
\]
\end{theorem}

Notice that this aligns with Theorem \ref{theo:Chain Rule for Two
Single} when $n = 2$. But we had one more chain rule!

\begin{theorem}{}{}
Let $f$ be a function of $n$ variables that are differentiable
functions of $m$ variables. Then, $f(x_1(t_1, \ldots, t_m), \ldots,
x_n(t_1, \ldots, t_m))$ is also differentiable and the following
chain rule holds.
\vspace{-0.6\baselineskip}
\[
\frac{\partial}{\partial t_i}
    f(x_1(t_1, \ldots, t_m), \ldots, x_n(t_1, \ldots, t_m))
= \sum_{j=1}^n
    \frac{\partial f}{\partial x_j} \frac{\partial x_j}{\partial t_i}
\]
\end{theorem}

Again for $n = 2$, this aligns with Theorem \ref{theo:Chain Rule for
Two Multi}. For proving this result, we can use the same one liner
to show it! Now let's move on to total differentials.

\begin{definition}{}{}
Let $f$ be a function of $n$ variables that is differentiable at
$(a_1, \ldots, a_n)$. The \textbf{total differential} of $f$ at the
point is defined as the following.
\vspace{-0.8\baselineskip}
\[
df = \sum_{i=1}^n f_{x_i}(a_1, \ldots, a_n) dx_i
\]
\end{definition}

For $n = 2$, this aligns with our Definition \ref{def:Total
Differential}. The intuition also transcribes here as well. We can
think $df$ as another function and that $\Delta f \approx df$ for
$dx_i \approx 0$. Continuing, let's generalize directional
derivatives.

\begin{definition}{}{}
Given a function $f$ of $n$ variables, the \textbf{directional
derivative} of $f$ at $(a_1, \ldots, a_n)$ in the direction of
a nonzero vector $\mathbf{v}$ is defined as the following.
\[
D_\mathbf{v} f(a_1, \ldots, a_n)
= \frac{1}{\| \mathbf{v} \|}
    \sum_{i=1}^n f_{x_i} (a_1, \ldots, a_n) v_i
\]
\end{definition}

Similarly, the \textbf{gradient} of $f$ at the given point is
defined as the following.
\[
\nabla f(a_1, \ldots, a_n)
= \langle
    f_{x_1}(a_1, \ldots, a_n), \ldots, f_{x_n}(a_1, \ldots, a_n)
\rangle
\]
Just like we observed before, we can also see that for any unit vector
$\mathbf{u}$, $D_\mathbf{u} f(a_1, \ldots, a_n) =
\nabla f(a_1, \ldots, a_n) \cdot \mathbf{u}$. Similarly, using
$\mathbf{v} \cdot \mathbf{w} = \| \mathbf{v} \| \| \mathbf{w} \|
\cos\theta$ again, we can see that the representation of gradients
and what its norm tells us also hold for functions of $n$ variables!

Now before we wrap up this note, let's generalize Theorem
\ref{theo:Level Curve Orthogonal}. There are generalizations of
this theorem for four or more variables, but for now let's focus
on functions of three variables.

\begin{theorem}{}{}
Let $f$ be a function of three variables that is differentiable at
$(x_0, y_0, z_0)$ and $\nabla f(x_0, y_0, z_0) \neq \mathbf{0}$.
Then, the tangent vectors of level surface through $(x_0, y_0, z_0)$
and $\nabla f(x_0, y_0, z_0)$ are orthogonal.
\end{theorem}
\begin{proof}
Let $S$ be the level surface $f(x, y, z) = k$ for some constant $k$.
Let $\mathbf{r}(t) = \langle x(t), y(t), z(t) \rangle$ be a
smooth curve on $S$ such that $\mathbf{r}'(t_0) = \mathbf{u}$ and
$(x(t_0), y(t_0), z(t_0)) = (x_0, y_0, z_0)$. Differentiating the
equation $f(x(t), y(t), z(t)) = k$ with respect to $t$,
\[
f_x(x_0, y_0, z_0) x'(t_0) +
f_y(x_0, y_0, z_0) y'(t_0) +
f_z(x_0, y_0, z_0) z'(t_0) = 0
\]
and $\nabla f(x_0, y_0, z_0) \cdot \mathbf{r}'(t_0) = 0$. Therefore,
the theorem holds.
\end{proof}

We can also see from the theorem that the equation of the tangent
plane can be represented as the following.
\[
f_x(x_0, y_0, z_0) (x - x_0) +
f_y(x_0, y_0, z_0) (y - y_0) +
f_z(x_0, y_0, z_0) (z - z_0) = 0
\]
This is it for this section! I think the fact that we can generalize
so many results from single variable calculus to functions of two
variables and that we can further generalize for functions of $n$
variables is what makes multivariable calculus so essential in many
different fields like physics and machine learning. In the next note,
we are going to apply the contents discussed here to find extrema
of some complicated functions.

\end{document}


